\section{Висновки:}

В ході виконання лабораторної роботи було успішно досягнуто поставлену мету: 

На основі виконаних досліджень та розробленого програмного комплексу сформульовано такі 
науково-практичні висновки:
\begin{enumerate}
    \item \textbf{Ознайомились з основами написання класів: } 
    Спроектовано та розподілено по окремих файлах базові класи сутностей (\texttt{Item}), 
    класи логічних структур (\texttt{StackCart}, \texttt{QueueCart}), сервісні компоненти 
    (\texttt{WarehouseManager}) та головний оркестратор (\texttt{Main}) згідно з вимогами
     \textit{Oracle Code Conventions}. \\
    \item \textbf{Освоїти інкапсуляцію класів:} 
    Всі внутрішні змінні станів (поля \texttt{price}, \texttt{name}, масиви 
    \texttt{stack} та \texttt{queue}) повністю ізольовані від зовнішнього втручання. 
    Доступ до них та валідація даних (наприклад, примусове скидання негативної ціни в 0) 
    реалізовано виключно через публічний інтерфейс методів та геттерів/сеттерів. \\
        
    \item \textbf{Розібралися із використанням конструктора:} 
    Для всіх класів розроблено кастомні конструктори з параметрами (ініціалізація об'єктів 
    товарів та виділення пам'яті під масиви заданої ємності). Для розв'язання конфліктів імен 
    між вхідними аргументами та полями класу у кучі (\textit{Heap}) успішно застосовано 
    ключове слово \texttt{this}. \\ 
    
\end{enumerate}


\newpage
\section{Інформаційні ресурси та посилання на вихідний код}

Згідно з вимогами до здачі лабораторної роботи, повний проект вихідного коду розробленого 
додатку розміщено у віддаленому репозиторії системи контролю версій. 

\begin{itemize}
    \item \textbf{Посилання на Git-репозиторій проєкту:} \\
    \texttt{https://github.com/nasuehfw/oop-lab2.git} % TODO: Тут моя реальная ссылка на github!
    \item \textbf{Електронний каталог спільного доступу:} \\
    \url{https://drive.google.com/drive/folders/1mIulRQSfoHAyQ2GnqdQbwGz2dFX-5VGK?usp=drive\_link}
\end{itemize}

